% ===================== Codificação, idioma e tipografia =====================
\usepackage[utf8]{inputenc}
\usepackage[T1]{fontenc}
\usepackage{lmodern}
\usepackage[brazilian]{babel}
\usepackage{csquotes}
\usepackage{microtype}
\usepackage[a4paper,margin=2.5cm]{geometry}
\usepackage{setspace}
\setstretch{1.1}
\setlength{\emergencystretch}{2em}

% ============================== Matemática ==================================
\usepackage{amsmath,amssymb,amsthm,mathtools}
\numberwithin{equation}{section}

% ====================== Tabelas, figuras, listas e caixas ===================
\usepackage{graphicx}
\usepackage{booktabs,array,tabularx,longtable,multirow}
\usepackage{xcolor}
\usepackage[font=small,labelfont=bf,labelsep=period]{caption}
\usepackage{float}
\usepackage{enumitem}
\setlist{itemsep=2pt,topsep=4pt}
\usepackage{tikz}
\usetikzlibrary{arrows.meta,positioning,calc,fit,shapes.geometric}
\usepackage[skins,breakable]{tcolorbox}

\newcolumntype{Y}{>{\centering\arraybackslash}X}
\newcolumntype{L}{>{\raggedright\arraybackslash}X}
\newcolumntype{R}{>{\raggedleft\arraybackslash}X}
\newcolumntype{P}[1]{>{\raggedright\arraybackslash}p{#1}}
\newtcolorbox{notabox}[1]{colback=yellow!8,colframe=orange!60!black,
  fonttitle=\bfseries,title={#1},boxrule=0.6pt,arc=2pt,breakable}

% =========================== Algoritmos e listagens =========================
\usepackage{algorithm}
\usepackage{algpseudocode}
\floatname{algorithm}{Algoritmo}
\algrenewcommand\algorithmicrequire{\textbf{Entrada:}}
\algrenewcommand\algorithmicensure{\textbf{Saída:}}
\algrenewcommand\algorithmicend{\textbf{fim}}
\algrenewcommand\algorithmicwhile{\textbf{enquanto}}
\algrenewcommand\algorithmicdo{\textbf{faça}}
\algrenewcommand\algorithmicfor{\textbf{para}}
\algrenewcommand\algorithmicforall{\textbf{para todo}}
\algrenewcommand\algorithmicif{\textbf{se}}
\algrenewcommand\algorithmicthen{\textbf{então}}
\algrenewcommand\algorithmicelse{\textbf{senão}}
\algrenewcommand\algorithmicreturn{\textbf{retorne}}
\algrenewcommand\algorithmicrepeat{\textbf{repita}}
\algrenewcommand\algorithmicuntil{\textbf{até que}}

\usepackage{listings}
\lstset{basicstyle=\ttfamily\footnotesize,breaklines=true,frame=single,
  columns=fullflexible,keepspaces=true,xleftmargin=2pt,xrightmargin=2pt}

% ============================ Teoremas e definições =========================
\theoremstyle{plain}
\newtheorem{proposicao}{Proposição}[section]
\theoremstyle{definition}
\newtheorem{definicao}[proposicao]{Definição}
\theoremstyle{remark}
\newtheorem{observacao}[proposicao]{Observação}
\addto\captionsbrazilian{\renewcommand{\proofname}{Demonstração}}

% ================================= Macros ===================================
\newcommand{\R}{\mathbb{R}}
\DeclareMathOperator{\sgn}{sgn}
\DeclareMathOperator*{\argmin}{arg\,min}
\DeclareMathOperator*{\argmax}{arg\,max}
\DeclareMathOperator{\lse}{lse}
\DeclareMathOperator{\softmax}{softmax}
\newcommand{\hamm}{d_{\mathrm{H}}}
\newcommand{\sstar}{s^{\ast}}
% Marcadores de conteúdo pendente (texto e células de tabela)
\newcommand{\apreencher}[1][]{\textcolor{red!70!black}{\textbf{[a preencher\ifx&#1&\else: #1\fi]}}}
\newcommand{\lacuna}{\textcolor{red!70!black}{$\cdots$}}

% ======================= Hiperlinks (carregar por último) ===================
\usepackage[colorlinks=true,linkcolor=blue!50!black,citecolor=blue!50!black,
  urlcolor=blue!50!black,unicode=true,
  pdftitle={Detecção de Anomalias em Telemetria de Sistemas Utilizando Redes de Hopfield},
  pdfauthor={Equipe}]{hyperref}
\usepackage{xurl}
